\documentclass[12pt]{article}
\usepackage[utf8]{inputenc}
\usepackage[T1]{fontenc}
\usepackage{hyperref}
\title{Gbagyi Language Resources for K-12 Curriculum Development}
\author{Compiled by Claude Sonnet 5.5 (claude-sonnet-5-5) for Andr\'as Kornai}
\date{}
\begin{document}
\maketitle

\section{Introduction}
Gbagyi (ISO 639-3: gbr; Glottocode: gbag1258) is an endangered Nupoid language of central Nigeria, with about 700,000 speakers. It is not known to be taught in schools. The search results show mainly a few grammars and one learner's book. No dictionary was found. The closely related Gbari variety is often treated separately, so check which variety a resource covers.

\section{Primary Educational Resources}
\begin{itemize}
\item \emph{Gbagyi Boda} (Shekwo, 2014), a learner's book for speakers who have lost fluency, covering the alphabet, pronunciation and thematic vocabulary: \url{https://nigeriareposit.nln.gov.ng/handle/20.500.14186/1489}
\end{itemize}

\section{Linguistic and Reference Materials}
\begin{itemize}
\item Shekwo (1979), \emph{Fundamentals of the Gbagyi language}, a grammar sketch (Glottolog record): \url{https://glottolog.org/resource/reference/id/109721}
\item Low (1908), \emph{Gbari grammar, notes and vocabulary} (Glottolog record; the Gbari variety): \url{https://glottolog.org/resource/reference/id/137614}
\item Numeral data for Gbagyi (MPG): \url{https://lingweb.eva.mpg.de/channumerals/Gbagyi.htm}
\item Wikipedia article on the Gwari language: \url{https://en.wikipedia.org/wiki/Gwari_language}
\end{itemize}

\section{Religious and Community Resources}
None were found in the search results.

\section{Digital and Audio Resources}
None were found in the search results.

\section{Curriculum Development Context}
A learner's book and two grammar sketches exist, but no dictionary, textbooks or digital resources were found. A curriculum would need substantial new development.
\end{document}